%%%%%%%%%%%%%%%%%%%%%%%%%%%%%%%%%%%%%%%%
%
% $ Autor: Gruppe 07 $
% $ Projekt: Magic Faucet Fountain $
% $ Pfad: LaTeXTemplate/Handbuch Magic Faucet Fountain/Chapters/de/Wartung.tex $
% $ Kurzbeschreibung: Dieses Dokument zeigt, wie der Magic Faucet Fountain gewartet werden soll. $
%
% !TeX spellcheck = de_DE/GB
% !TeX program = pdflatex
% !BIB program = biber/bibtex
% !TeX encoding = utf8
%
%%%%%%%%%%%%%%%%%%%%%%%%%%%%%%%%%%%%%%%%

\chapter{Wartung}
	\textbf{Allgemeiner Hinweis:}\\
	Für jegliche Wartungsarbeiten muss die Stromzufuhr unterbrochen werden. Die Stromzufuhr darf erst nach Beendigung der Wartungsarbeiten eingeschaltet werden 

\textbf{Wartungshinweise:}
	\begin{enumerate}
		\item Nach längeren Standzeiten sollte das Wasser im Blumentopf gewechselt werden. Hierfür wird der Deckel mit dem montierten Acrylrohr vorsichtig abgenommen. Das alte Wasser wird entleert und der Blumentopf bei Bedarf gründlich gereinigt. Anschließend wird frisches Wasser bis zur markierten Füllhöhe eingefüllt, und der Deckel wieder aufgesetzt.
		\item Das Schutzrohr, in dem sich der Ultraschallsensor befindet, ist fest verbaut und darf nicht entfern t werden. Zur Reinigung des Sensors genügt es, den Deckel abzunehmen und den Sensor mit einem trockenen, fusselfreien Mikrofasertuch vorsichtig abzuwischen. Direkter Hautkontakt mit der Sensorfläche ist zu vermeiden.
		\item Die Pumpe sollte regelmäßig auf Verschmutzungen überprüft und gereinigt werden. Dazu wird der Deckel mit dem Acrylrohr abgenommen und die Pumpe vorsichtig aus ihrer Halterung gelöst. Anschließend wird sie unter klarem Wasser ausgespült. Nach der Reinigung wird die Pumpe wieder eingesetzt und der Deckel ordnungsgemäß aufgesetzt.
	\end{enumerate}
	 


